\chapter{Remediation Report}

\section{Overview}

This chapter provides recommended remediation measures for the vulnerabilities identified during the penetration test. 
\\

It outlines the actions required to mitigate or eliminate the identified security weaknesses, reduce the likelihood of successful exploitation, and strengthen the overall security posture of the target system. 
\\

The recommendations are prioritised according to the severity and potential impact of the associated findings.

\section{Remediations}

\subsection{Immediate Mitigation}

\subsubsection{SUID Binary Search Path (F-006)}
\begin{itemize}
    \item \textbf{Issue:} \texttt{easysysinfo} executes \texttt{cat} via \texttt{system()} without absolute paths, allowing local root escalation via \texttt{PATH} hijacking.
    \item \textbf{Remediation:}
    \begin{enumerate}
        \item Modify source code to specify absolute execution paths (e.g., \texttt{/bin/cat}) or replace shell execution with direct system calls like \texttt{execve()}.
    \end{enumerate}
\end{itemize}

\subsubsection{Web Application Path Traversal and LFI (F-002, F-003)}
\begin{itemize}
    \item \textbf{Issue:} Unsanitized input in \texttt{randylogs.php} allows reading arbitrary files (e.g., \texttt{/var/log/auth.log}), which can be paired with SFTP log poisoning for Remote Code Execution.
    \item \textbf{Remediation:}
    \begin{enumerate}
        \item \textbf{Fix LFI:} Remove direct passing of user input into \texttt{include}, \texttt{require}, or file-reading functions. Implement a strict array-based whitelist mapping user input keys to static, pre-approved file paths.
        \item \textbf{Restrict Access:} Adjust file permissions so system logs are not readable by the web server user (\texttt{www-data}).
    \end{enumerate}
\end{itemize}

\subsubsection{Sensitive Files and Backup Artifacts (F-001, F-004, F-005)}
\begin{itemize}
    \item \textbf{Issue:} Publicly accessible operational files, weak backup archive passwords, and insecure permissions on \texttt{/var/backups/user\_backups.zip}.
    \item \textbf{Remediation:}
    \begin{enumerate}
        \item Relocate sensitive files (\texttt{tasks\_todo.txt}, logs, developer notes) outside the public web root directory.
        \item Disable directory browsing across web server configurations.
        \item Restrict file permissions on backup archives so only the owner can access them.
        \item Re-encrypt backup archives using AES-256 encryption with a high-entropy passphrase.
    \end{enumerate}
\end{itemize}

\subsection{Software Upgrades}

\subsubsection{Upgrade Apache HTTP Server (F-007 to F-084)}
\begin{itemize}
    \item \textbf{Scope:} Addresses 78 reported CVEs (including request smuggling, heap buffer overflows, HTTP/2 DoS, and SSRF flaws).
    \item \textbf{Remediation:}
    \begin{enumerate}
        \item Upgrade the Apache HTTP Server package to version 2.4.68 or later.
    \end{enumerate}
\end{itemize}

\subsubsection{Operating System Kernel Patching (F-089)}
\begin{itemize}
    \item \textbf{Issue:} Linux kernel versions 4.15.x to 4.19.x mishandle user namespace mapping (CVE-2018-18955), allowing local privilege escalation.
    \item \textbf{Remediation:}
    \begin{enumerate}
        \item Update the Linux Kernel package to the latest distribution release.
    \end{enumerate}
\end{itemize}